\subsection{Objetivos de negocio}

\subsubsection{Propósito y visión global}

Aprovechando la infraestructura tecnológica del CAETEC, transformar la gestión del establo lechero hacia un modelo de ganadería de precisión 4.0 proactiva. El proyecto fusiona la telemetría continua de los 3 robots con la captura multimodal de dispositivos móviles accesibles (iPhone RGB + LiDAR), transformando datos de sensores en decisiones operativas automatizadas que preserven el bienestar animal, eviten mermas productivas por estrés y aseguren la viabilidad reproductiva de las 160 vacas Holstein del hato.

\subsubsection{Objetivo general del negocio}

Permitir evaluar de forma frecuente y no invasiva la Condición Corporal (BCS) de cada vaca, con equipo accesible para el personal del establo, buscando prevenir caídas de BCS por debajo de 2.75 para evitar el colapso de fertilidad postparto. \textbf{Objetivo secundario:} anticipar y mitigar la saturación operativa de las 3 estaciones DeLaval VMS del CAETEC, para reducir las esperas en la manga, mejorar la distribución del tráfico de las vacas a lo largo del día y prevenir caídas de producción por eyección retardada. Todo esto orientado a elevar los litros de leche recolectados, garantizar el bienestar de los animales y optimizar la rentabilidad del socio.

\subsubsection{Objetivos específicos del negocio por eje estratégico}

\paragraph{Eje A (principal): salud, bienestar y condición corporal (BCS)}

\begin{itemize}
  \item \textbf{Evaluación objetiva y frecuente de la condición corporal:} sustituir la evaluación visual esporádica del BCS por una evaluación frecuente, no invasiva y consistente, realizada en el propio establo, que permita dar seguimiento a cada vaca a lo largo de la lactancia.
  \item \textbf{Viabilidad reproductiva:} detectar a tiempo a las vacas en postparto cuyo BCS tiende a caer por debajo de 2.75, para ajustar su nutrición y manejo antes de que se comprometa la preñez. (Se apoya en el artículo sobre BCS y fertilidad revisado en 1.1.1).
  \item \textbf{Prevención proactiva de patologías metabólicas:} alertar a los nutricionistas y veterinarios cuando una vaca entre en periodo seco con sobrepeso (BCS > 4) o presente emaciación severa (BCS < 2.25), previniendo cetosis clínica, hígado graso e hipocalcemia.
\end{itemize}

\paragraph{Eje B: eficiencia operativa, capacidad y teoría de colas en el VMS}

\begin{itemize}
  \item \textbf{Mitigación de la eyección retardada de leche (DME) y del flujo bimodal:} con base en la literatura (artículo sobre eyección retardada revisado en 1.1.1, donde el flujo bimodal genera pérdidas de 1.2 kg de leche por sesión), reducir los tiempos muertos y el nerviosismo en la manga de acceso para asegurar que la estimulación del pezón induzca una bajada continua y homogénea de leche.
  \item \textbf{Protección de vacas primíparas:} evitar que las vacas dominantes de alta producción lleguen a monopolizar el VMS, garantizando que otras vacas alcancen la frecuencia óptima de ordeño sin sufrir desplazamientos agresivos en la fila.
  \item \textbf{Anticipación de colapsos:} evitar que cualquiera de los 3 robots opere saturado durante periodos prolongados, dando al personal de campo aviso con anticipación suficiente para rebalancear el tráfico de las 160 vacas antes de que se forme una cola física prolongada.
\end{itemize}

\subsubsection{Indicadores clave de desempeño del negocio (KPI)}

\noindent
\begin{tabularx}{\textwidth}{>{\raggedright\arraybackslash}X>{\raggedright\arraybackslash}X>{\raggedright\arraybackslash}X>{\raggedright\arraybackslash}X}
  \toprule
  Indicador de negocio & Línea base & Meta esperada & Impacto cuantitativo en el establo del CAETEC \\
  \midrule
  \textbf{Pérdida por flujo bimodal (DME)} & Merma recurrente no monitoreada de 1.2 kg por ordeño atípico. & Reducción del 40\% en sesiones bimodales por estrés de espera. & Recuperación de hasta 150--200 kg de leche diarios en el hato de 160 vacas. \\
  \addlinespace
  \textbf{Horas diarias de saturación robótica (>90\%)} & 4--6 horas/día con colapso silencioso de manga. & Reducción del 35\% de horas de colapso mediante alertas anticipadas. & Aprovechamiento óptimo de los 3 robots DeLaval VMS del CAETEC. \\
  \addlinespace
  \textbf{Tiempo de espera promedio en la manga} & Picos de hasta 45--90 min en horas pico de alimento. & $\leq$20 minutos sostenidos en horas pico. & Mayor tiempo de rumia y aumento estimado de 1.0 a 1.8 kg de leche/vaca/día. \\
  \bottomrule
\end{tabularx}
